\documentclass[10pt,a4paper]{article}
\usepackage[margin=1.8cm]{geometry}
\usepackage{booktabs,tabularx,enumitem,hyperref,xcolor,amsmath,array}
\usepackage[T1]{fontenc}
\usepackage{lmodern}
\hypersetup{colorlinks=true,linkcolor=black,urlcolor=blue!50!black}
\setlist{nosep,leftmargin=*}
\setlength{\parskip}{3pt}\setlength{\parindent}{0pt}
\newcommand{\code}[1]{\texttt{\small #1}}
\newcommand{\meas}{\textsuperscript{\textbf{M}}}
\newcommand{\assu}{\textsuperscript{\textbf{A}}}
\renewcommand{\arraystretch}{1.12}

\title{\vspace{-1.2cm}Knowledge Distillation and Deployment of a Command Parser on a Local Cloud\\
\large CS3081 Cloud Computing --- Mini Project Report}
\author{Manas Trivedi \quad (230101073)\\
\small Repository: \url{https://github.com/Manas-Trivedi/cc-kd-command-parser} \quad (assessed commit: tag \code{final})}
\date{}

\begin{document}\sloppy
\maketitle\vspace{-1cm}

\paragraph{Overview.} A teacher LLM (phi3.5, 3.8\,B, on Ollama) labels 3,000 voice-assistant commands with a
four-key JSON frame. A 70\,M-parameter student (Pythia-70M) is fine-tuned on those replies
(sequence-level distillation) and served by a FastAPI service on a k3d cluster. Weights live in
SeaweedFS object storage and are downloaded by every pod at start-up. All evidence is in \code{results/}.
\textbf{Notation:} values marked \meas{} are measured on my laptop (Apple M3 Pro, 11 cores, 18\,GB; Docker via
OrbStack with 9\,GB); values marked \assu{} are assumptions or given prices.

{\scriptsize
\begin{verbatim}
 laptop: Ollama phi3.5 (Apple GPU)  <-- HTTP --  ns pipeline: Job teacher-label (Indexed, 3 pods, SHARD_INDEX = index)
                                                   | raw replies -> s3://datasets/teacher/shard-{0,1,2}.jsonl
 ns data: SeaweedFS (S3 API, PVC)  datasets/{commands_pool.jsonl, teacher/, v1/train.jsonl}   models/{v1/, v2/}
                                                   | boto3 download at pod start (keys from Secret s3-credentials)
 localhost:8080 -> Traefik Ingress /parse -> Service app/parser -> Deployment parser: 2 pods (HPA 1-2 @ 70% CPU)
\end{verbatim}}

%==================================================================
\section{Task 1 --- Object storage and parallel teacher labelling}
\textbf{Setup.} \code{cluster/k3d.yaml} defines one server and one agent node, maps the Traefik load balancer to
\code{localhost:8080} and the SeaweedFS S3 NodePort to \code{localhost:8333}. SeaweedFS runs as one pod
(\code{weed server -s3}) on a PersistentVolumeClaim in namespace \code{data}. Its S3 identity file and the
client keys are generated by \code{scripts/secrets.sh} into \code{.secrets/} (git-ignored) and loaded as
Kubernetes Secrets, so no credential is in the repository. Buckets \code{datasets} and \code{models} exist and
\code{datasets/commands\_pool.jsonl} is uploaded (seed Job, or \code{scripts/s3.py}).

\textbf{Teacher.} Ollama runs natively on the laptop so phi3.5 uses the Apple GPU; pods reach it at
\code{host.docker.internal:11434}. Options: \code{temperature 0, seed 42, num\_predict 96}, and Ollama's
\code{format: "json"}. The prompt (\code{prompts/teacher\_prompt.txt}) lists the eight intents, the three slot
rules and ten examples I wrote myself (none from the validation set). In a pilot on the first 80 pool commands,
phi3.5 without JSON mode appended explanatory ``Note:'' paragraphs after the JSON in 13/80 replies\meas{}; with JSON
mode 0/80 were invalid\meas{}, so I kept JSON mode. Replies are stored exactly as generated.

\textbf{Sharding (worker/worker.py).} Each worker reads \code{SHARD\_INDEX} and \code{SHARD\_COUNT}, keeps the
pool lines whose 0-based line number $i$ satisfies $i \bmod \text{SHARD\_COUNT} = \text{SHARD\_INDEX}$, and writes
\code{\{id, text, teacher\_raw\}} lines to its own object \code{datasets/teacher/shard-<k>.jsonl}. The workers run
as one Kubernetes \emph{Indexed Job} (\code{completions = parallelism = 3}) in namespace \code{pipeline}; Kubernetes
gives each pod a distinct completion index $0..2$, which the manifest passes in as \code{SHARD\_INDEX}. Modulo
(round-robin) sharding gives equal shard sizes (1,000 each) and mixes intents evenly, and because each line maps to
exactly one remainder there can be no gap or overlap. Each worker re-uploads its shard every 20 commands and, on
start, skips ids already in its shard object. This mattered: the Docker VM froze mid-run and had to be restarted;
the Job re-created the three pods, and each resumed at 300/1,000 instead of starting over
(\code{results/task1/label\_resume\_after\_restart.txt}). Total wall time 76\,min\meas{} (3 parallel requests to
one Ollama, $\approx$0.55 commands/s per worker\meas{}).
\code{results/task1/shard\_coverage.txt} confirms all 3,000 ids appear exactly once\meas{}.

\textbf{Cleaning (worker/clean.py).} Records are processed in id order; the first failed check is the reason.
Survivors are de-duplicated on normalised text (lower case, single spaces), keeping the smallest id. Output is
serialised deterministically (sorted keys, fixed separators), so re-running gives the same SHA-256; I re-ran it on
the downloaded shards and got the same checksum\meas{}.

\begin{center}\small
\begin{tabular}{lrrrrrrr}
\toprule
& Input & (a) not JSON & (b) wrong keys & (c) bad intent & (d) slot not in command & (e) duplicates & Kept\\
\midrule
Records\meas{} & 3,000 & 2 & 2 & 13 & 162 & 49 & 2,772\\
\bottomrule
\end{tabular}\\[2pt]
$2+2+13+162+49+2772 = 3000$. \quad \code{train.sha256} = \code{a58e34e5f831aea4\ldots c9ec5e} (full value in \code{data/v1/train.sha256}).
\end{center}
Most rejections are (d): the teacher paraphrases (``in two days'' for ``after two days'') or invents
(\code{"time": "now"}), which the copy-exactly rule forbids. 50 commands in the pool are repeats; one repeat was
already rejected by (d), so (e) removes 49.

%==================================================================
\section{Task 2 --- Training the student (sequence-level distillation)}
\textbf{Method.} Each example is \code{Command: <text>\textbackslash nJSON: <teacher JSON><eos>}; the loss is
masked to the reply tokens, so the student learns $p(\text{teacher reply}\mid\text{command})$. Training uses a plain
PyTorch loop (AdamW, linear warm-up/decay, gradient clip 1.0) on the CPU, on the first 1,500 cleaned records.
\code{training/config.yaml} holds seed 42, learning rate $5\times10^{-4}$, batch 16, 10 epochs, max length 96 and
the library versions (torch 2.5.1, transformers 4.46.3, \ldots); \code{make train} reads only this file and
\code{make train SUBSET=50} trains on 50 records. Seeds (Python, NumPy, torch, data order) are set in code and
deterministic algorithms are on: re-training v1 reproduced byte-identical predictions\meas{}.
One epoch takes $\approx$31\,s\meas{} on 11 CPU threads. Inference is greedy and stops at EOS or the first newline.

\textbf{Versions.} My first run used 8 epochs; it left one of 100 validation outputs as truncated JSON
(99\% validity\meas{}, kept in \code{results/experiments/}). Because each invalid output becomes an HTTP 500 in the
load test, I increased to 10 epochs: this is \textbf{v1} (seed 42). \textbf{v2} is identical except seed 7.
Training loss (v1, \code{training/loss.csv}\meas{}): 0.808, 0.392, 0.219, 0.157, 0.107, 0.084, 0.049, 0.027, 0.012, 0.004.
Weights are saved as float16 (141\,MB instead of 282\,MB; my upload link managed only $\approx$40\,KB/s\meas{},
and the first 226\,MB Release upload timed out) and loaded back as float32 for inference; validation predictions
are byte-identical to the float32 weights\meas{}. Running training twice from scratch gave the same weights file
SHA-256 and the same losses\meas{} (\code{results/retrain\_reproducibility.txt}).
Model and tokenizer are in \code{s3://models/v1/} and \code{s3://models/v2/}. Each version is packed into a
reproducible tarball (sorted entries, zero mtimes); \code{models/v1/model.sha256} is the SHA-256 of that tarball,
which is the GitHub Release asset \code{model-v1.tar.gz} (\code{18a4f792\ldots 8731571a}); CI re-checks it with
\code{sha256sum -c}.

\begin{center}\small
\begin{tabular}{lrrrrrr}
\toprule
Validation (100 commands)\meas{} & JSON valid & Intent & Date & Time & Place & Full frame\\
\midrule
Teacher phi3.5 (same prompt/settings) & 99 & 90 & 77 & 84 & 91 & \textbf{54}\\
Untrained Pythia-70M & 0 & 0 & 0 & 0 & 0 & 0\\
Student v1 (10 epochs, seed 42) & \textbf{100} & 86 & 78 & 85 & 82 & \textbf{45}\\
Student v2 (10 epochs, seed 7) & 100 & 86 & 80 & 81 & 85 & 48\\
\emph{First attempt (8 epochs, seed 42)} & 99 & 83 & 72 & 83 & 81 & 40\\
\bottomrule
\end{tabular}
\end{center}
\textbf{What changed.} The untrained model just continues text (\code{- name: "calendar"}, \code{\{}) and never
produces a valid frame. After distillation the student always emits the schema and reaches 83\% of the teacher's
full-frame accuracy (45 vs 54) with 1/54 of the parameters; it agrees exactly with the teacher on 52/100
commands\meas{}. It loses most on intent (86 vs 90) and place (82 vs 91): place names are rare tokens it has seen
few times. The student is \emph{more} valid than the teacher: the teacher's one invalid output was
\code{"intent": null}, a mistake the student, which only ever saw cleaned replies, never learned. Raw outputs and
score reports are in \code{results/} (\code{teacher\_validation.json}, \code{untrained\_validation.json},
\code{student\_v\{1,2\}\_validation.json} and the matching \code{*\_predictions.jsonl}).

\textbf{Sequence-level vs logit-based distillation.}
\emph{Sequence-level} KD (Kim \& Rush, 2016) trains the student with ordinary cross-entropy on the
\emph{text the teacher generated} (here the teacher's greedy JSON reply): a single hard target sequence per input.
\emph{Logit-based} (word-level) KD trains the student to match the teacher's full \emph{probability distribution
over the vocabulary at every position} (e.g.\ KL divergence between softened teacher and student softmax outputs,
as in Hinton et al.\ and TRL's GKD/DistillationTrainer), so it also learns how likely the alternatives were.
\emph{It is not possible here} because (i) the teacher is reached only through Ollama's text API, which returns
generated text, not per-token logits over the vocabulary; and (ii) phi3.5 and Pythia use different tokenizers
(32k vs 50k vocabularies), so their distributions are over different token sets and cannot be compared position by
position without an alignment scheme. \emph{Limitation of my approach:} the student only sees one hard answer per
command, including the teacher's mistakes: the teacher scores 54\% full-frame, so its errors (e.g.\ labelling
``remind me'' as \code{alarm\_set}, or extra date words) are copied as if correct. The student's ceiling is the
teacher, and a hard label carries no information about how uncertain the teacher was. Cleaning removes only
\emph{format} errors; wrong-but-valid labels stay in the training set.

\textbf{Suitability of Pythia-70M.} Apache-2.0 licensed and small enough to train in about 5 minutes on a laptop CPU, which makes it good for this
demonstration. But it is a base model pre-trained on the Pile with no instruction tuning, and the model card states it
is not intended for deployment; with 70\,M parameters it memorises the training distribution (loss
$\to 0.004$) and generalises poorly to unseen place names. It is adequate for a fixed-format task, not for
open-ended use.

%==================================================================
\section{Task 3 --- Inference service in a container}
\code{service/app.py} (FastAPI, one uvicorn worker) exposes \code{POST /parse}, \code{GET /healthz} (always 200
while the process runs) and \code{GET /readyz} (503 until the model is loaded, then 200). Configuration comes from
\code{MODEL\_URI}, \code{MODEL\_VERSION}, \code{S3\_ENDPOINT}, \code{MAX\_NEW\_TOKENS} and \code{FAULT\_INJECT}; the
S3 keys come from the Secret \code{s3-credentials} via \code{envFrom}. At start-up a background thread lists
\code{MODEL\_URI}, downloads every object with boto3, loads the model and runs one warm-up generation. Every
output is parsed and validated with \code{jsonschema} against \code{schema.json} (exact keys, intent enum, string or
null slots); if invalid, the response is HTTP 500 \code{\{"error": "invalid\_model\_output"\}}. With
\code{FAULT\_INJECT=1} the output is replaced with \code{not json} before validation. Successful responses carry
\code{X-Model-Version} and \code{X-Served-By} (pod name from the Downward API). Inference is serialised with a lock
and uses one torch thread to match the 1-CPU limit. Nine unit tests cover the endpoints, headers, six kinds of
invalid output and fault injection.

\textbf{Image.} \code{python:3.12.8-slim} + the CPU-only torch wheel from \code{download.pytorch.org/whl/cpu}
(the default wheel pulls CUDA libraries) + a fully pinned \code{service/requirements.txt} (46 packages incl.\ torch, generated
with \code{uv pip compile}). No weights are copied in; a \code{find} for \code{*.safetensors}/\code{*.bin} in the
image returns 0\meas{}. Size \textbf{1.23\,GB}\meas{}. Build: \code{docker build -f service/Dockerfile .} (or
\code{make images}).

\textbf{Running under limits} (\code{scripts/container\_test.sh}, \code{--cpus 1 --memory 1g}; output in
\code{results/task3/container\_test.txt}). 50 sequential validation commands: 50/50 HTTP 200, median latency
408\,ms, max 468\,ms\meas{}. \textbf{Peak memory: 788.7\,MiB} process peak RSS (\code{VmHWM})\meas{}; the
cgroup's \code{memory.peak} was 765.5\,MiB and the highest \code{docker stats} sample 593.1\,MiB\meas{}, all under
the 1,024\,MiB limit. The peak happens while the float16 file is converted to float32 at load; steady state is
$\approx$590\,MiB. No OOM kill, 0 restarts. \code{FAULT\_INJECT=1} returned
\code{HTTP 500 \{"error":"invalid\_model\_output"\}}\meas{}.
\emph{Problem found:} the first version took 9.8\,s per request at 1 CPU\meas{}. Profiling showed the oneDNN
(mkldnn) kernels in the arm64 Linux torch wheel are pathologically slow for batch-1 decoding; disabling them
(\code{torch.backends.mkldnn.enabled = False}) brought it to 0.43\,s\meas{}, a 23$\times$ speed-up.

%==================================================================
\section{Task 4 --- Kubernetes deployment, scaling and reliability}
\code{k8s/app/parser.yaml}: namespace \code{app}; Deployment with 2 replicas, requests 500m CPU / 512\,Mi, limits
1 CPU / 1\,Gi; liveness probe on \code{/healthz} (every 10\,s, 3 failures) and readiness probe on \code{/readyz}
(every 2\,s, 1 failure); rolling update \code{maxUnavailable: 0, maxSurge: 1}; a \code{preStop} sleep of 8\,s so a
pod keeps serving requests already routed to it while it is removed from the endpoints; ClusterIP Service; Traefik
Ingress routing \code{/parse}; HPA (\code{autoscaling/v2}) min 1, max 2, 70\% average CPU, scale-down
stabilisation 60\,s. Experiments are scripted in \code{scripts/k8s\_tests.sh}; outputs in \code{results/task4/}. (The HPA, rolling-update, recovery and load-balancing runs were done before I switched the stored weights to
float16; they served the float32 copy of the same v1/v2 models, which gives identical outputs. Cold start and
Task 7 were re-run on the float16 copy.)

\begin{center}\small
\begin{tabularx}{\textwidth}{>{\raggedright}p{2.7cm}X}
\toprule
Test & Result\meas{}\\
\midrule
Ingress & \code{curl -X POST localhost:8080/parse} returns the frame with both headers (\code{resources.txt}).\\
No endpoints while loading & Scaled 0$\to$1 and polled each second: pod \code{/readyz}=503, EndpointSlice
\code{ready=false}, Ingress \code{/parse} = HTTP 503, until the log line ``model v1 ready'', then endpoint
\code{true} and 200 (\code{probes.txt}).\\
HPA & Idle at 2 replicas, 0\% CPU $\Rightarrow$ scaled to 1 (05:06:40). Under the 4-user load the single pod hit
990m (198\% of its request) $\Rightarrow$ \emph{New size: 2} at 05:07:21. After the load ended CPU fell to 0\% and
the HPA scaled back to 1 (event ``All metrics below target'', 05:11:49) (\code{hpa.txt}).\\
Rolling update & During a 5-min load test: \code{kubectl set env} to \code{MODEL\_VERSION=v2},
\code{MODEL\_URI=s3://models/v2/}, then \code{kubectl rollout undo}. \textbf{1,314 requests, 0 failures (0.00\%)};
v2 served 307 of them; after the rollback 5/5 probes returned \code{X-Model-Version: v1} (\code{rollout.txt}).\\
Recovery & Deleted one of two pods under load: Ready count went 2$\to$1$\to$2 in 7\,s, the old pod was gone after
11\,s; \textbf{774 requests, 0 failures} (\code{recovery.txt}).\\
Cold start & Deleted all pods 3 times. Pod creation$\to$Ready (slowest pod per run): 4, 4, 4\,s,
\textbf{median 4\,s} (1\,s timestamp resolution). Logs show \code{downloading model s3://models/v1/ \ldots
download finished: 144.4 MB in 0.2\,s \ldots model v1 ready} (\code{coldstart.txt}).\\
Load balancing & Exactly two Ready replicas (HPA min pinned to 2 for the run), standard command through the
Ingress: \textbf{253 (49.8\%) and 255 (50.2\%)} of 508 successful requests (\code{balance.txt}).\\
\bottomrule
\end{tabularx}
\end{center}
\textbf{Limitations of the local cluster.} Both k3d ``nodes'' are containers on one laptop: losing the laptop, the
Docker VM (which actually froze once) or the disk loses everything, so this shows \emph{pod}-level recovery, not
fault tolerance across machines or zones. SeaweedFS is a single replica on a local-path volume: a single point of
failure. metrics-server samples every 15\,s and the HPA syncs every 15\,s, so scale-up takes about 40--60\,s;
\code{max 2} replicas is the ceiling. CPU ``utilisation'' is relative to the 500m request, so a pod at its 1-CPU limit shows
200\%. Pods are throttled by the shared host CPU, so other load on the laptop changes the results.

%==================================================================
\section{Task 5 --- Subnetting and network access control}
Largest first from 10.20.0.0, smallest prefix with $2^{32-p}-2 \ge$ hosts (working in \code{network/subnets.md}):

\begin{center}\small
\begin{tabular}{lrcllr}
\toprule
Subnet & Needed & Prefix & Network & Usable range / Broadcast & Usable\\
\midrule
app & 1,000 & /22 & 10.20.0.0 & 10.20.0.1--10.20.3.254 / 10.20.3.255 & 1,022\\
public & 250 & /24 & 10.20.4.0 & 10.20.4.1--10.20.4.254 / 10.20.4.255 & 254\\
data & 100 & /25 & 10.20.5.0 & 10.20.5.1--10.20.5.126 / 10.20.5.127 & 126\\
management & 20 & /27 & 10.20.5.128 & 10.20.5.129--10.20.5.158 / 10.20.5.159 & 30\\
\bottomrule
\end{tabular}
\end{center}
\textbf{NACL for the public subnet} (\code{network/nacl.json}). Inbound: 90 deny \emph{all} from
198.51.100.0/24; 100 allow TCP 80 and 110 allow TCP 443 from 0.0.0.0/0; 120 allow TCP 22 from 10.20.5.128/27;
130 allow TCP 1024--65535 from 10.20.0.0/22 (replies from app to the load balancer's ephemeral ports); everything
else falls to the default deny. Outbound: 100 allow TCP 8000 to 10.20.0.0/22; 110 allow TCP 1024--65535 to
0.0.0.0/0 (replies to clients). \code{nacl\_sim.py --check}: \textbf{6/6 correct}\meas{}.

\textbf{Why replies need rules, and why order matters.} A NACL is \emph{stateless}: it keeps no record of
connections, so a reply packet is checked against the rules like any new packet. The reply to a client's request
to port 443 leaves for the client's ephemeral port (1024--65535) and needs outbound rule 110; the app's reply to
the load balancer's call to port 8000 returns to the load balancer's ephemeral port and needs inbound rule 130.
(A stateful security group would allow these automatically.) Rules are evaluated from the lowest number and the
\emph{first match wins}. A packet from 198.51.100.9 to port 443 matches both ``deny 198.51.100.0/24'' and
``allow 443 from anywhere''; it is blocked only because the deny is rule 90, lower than the allow rules 100/110. Numbered,
say, 200, the abusive range would get in on the web ports.

%==================================================================
\section{Task 6 --- Reproducibility, testing and CI/CD}
\textbf{Fresh clone in three commands:} \code{git clone \ldots}, \code{cd cc-kd-command-parser}, \code{make up}.
\code{make up} creates the cluster from \code{cluster/k3d.yaml}, generates credentials into Secrets, deploys SeaweedFS, builds both
images and imports them into k3d, runs a seed Job that creates the buckets, uploads the pool and copies the
v1/v2 Release tarballs (after checking their SHA-256) into \code{s3://models/}, and deploys the app.
\code{make down} deletes the cluster. I tested this after \code{make down}, from a fresh
clone into an empty directory: \code{make up} exited 0 and \code{/parse} answered through the Ingress with
\code{X-Model-Version: v1} (\code{results/task6/make\_up\_fresh\_clone.log})\meas{}. It took 22.5\,min\meas{}, almost all
of it downloading the k3s system images and the two 127\,MB Release tarballs over my $\approx$0.5--1\,MB/s link. A first
attempt timed out waiting for SeaweedFS while k3s was still pulling its system images, so I raised the timeouts.
Prerequisites: Docker, k3d, kubectl, make, openssl (and Python 3 for the training/test targets).

\textbf{GitHub Actions} (\code{.github/workflows/ci.yml}), stages in order: \emph{lint} (ruff) $\to$
\emph{unit-tests} (pytest, 13 tests: sharding, the cleaning rules and determinism, reproducible packaging,
service endpoints/headers/validation/fault injection) $\to$ \emph{evaluate} (download the Release with
\code{gh release download}, verify the checksum, run the student on \code{validation.jsonl}, \code{score.py})
$\to$ \emph{quality-gate} (\code{score.py --min-validity 95}); plus a parallel \emph{secret-scan} job (gitleaks over
the full history). A manual input \code{model=untrained} evaluates the base Pythia-70M instead.
\textbf{Runs:} passing run on the student v1 Release (checksum \code{OK}, validity 100, full frame 45; same
numbers on GitHub's x86 runner as on my laptop):
\url{https://github.com/Manas-Trivedi/cc-kd-command-parser/actions/runs/42317803924}.
\textbf{Gate doing its job} (untrained Pythia-70M, \code{QUALITY GATE FAILED: json\_validity 0.0 < 95.0}; lint,
tests and evaluation passed, the quality-gate job failed):
\url{https://github.com/Manas-Trivedi/cc-kd-command-parser/actions/runs/42318098565}.
gitleaks scans the full git history and reports no leaks (CI job and \code{results/task6/gitleaks\_local.txt}); the
only secrets on disk are in the git-ignored \code{.secrets/}. Evidence: \code{results/task6/}.

%==================================================================
\section{Task 7 --- Performance, cost and review}
Standard command (4 users, 2 min, through the Ingress) against the deployment in its normal state (HPA at
1 replica when idle). Files: \code{results/task7/}.

\begin{center}\small
\begin{tabularx}{\textwidth}{>{\raggedright}p{3.0cm}X>{\raggedright\arraybackslash}p{3.3cm}}
\toprule
Metric & Value\meas{} & Source\\
\midrule
Requests / failures & 467 / 0 (\textbf{error rate 0.00\%}) & \code{loadtest\_stats.csv}\\
Latency p50 / p95 & \textbf{1,300\,ms / 1,700\,ms} (min 405, max 1,880) & \code{loadtest\_stats.csv}\\
Throughput & \textbf{3.91 req/s} (4.3--4.4 req/s once 2 replicas were up) & \code{loadtest\_stats.csv}\\
Pod CPU / memory & pod 1: 980--991m, 583--620\,Mi; pod 2 (added by the HPA during the run): 801--998m, 588--613\,Mi & \code{perf\_top.txt}\\
Cold start & median 4\,s creation$\to$Ready (Task 4) & \code{task4/coldstart.txt}\\
\bottomrule
\end{tabularx}
\end{center}
Both pods run at their 1-CPU limit: the service is CPU-bound. One request takes $\approx$0.4\,s of CPU time
($\approx$25 greedy decode steps), so 4 concurrent users queue behind the inference lock, which explains a
1.3\,s median. The replica split was 267/202 because the second pod only started partway through the run.

\textbf{Cost} (rate card\assu{}: 3.50 INR per vCPU-hour, 0.45 INR per GiB-hour; billed on limits).
Per replica: $1\text{ vCPU}\times3.50 + 1\text{ GiB}\times0.45 = 3.95$ INR/h. Under sustained load the HPA runs
2 replicas: $2\times3.95 = 7.90$ INR/h. Requests per hour at the measured 3.9116 req/s:
$3.9116\times3600 = 14{,}082$. Cost per 1,000 requests $= 7.90 / 14{,}082 \times 1000 = \textbf{0.561 INR}$.
(At the 2-replica steady-state throughput of 4.39 req/s\meas{}: $7.90/15{,}804\times1000 = 0.50$ INR.)
The cost is dominated by the CPU limit (89\% of the hourly price).

\textbf{SLO.} \emph{99\% of \code{/parse} requests return HTTP 200 within 30\,s, measured over one standard
load-test run.} Window: one 2-minute run (467 requests). Error budget $= 1\% \times 467 = 4.67$, i.e.\ \textbf{4
failed requests} (rounded down). Used: 0 failures = 0\% of the budget\meas{}. (Rolling update and recovery runs also used 0 of
their budgets of 13 and 7 requests.)

\textbf{Review.}\\[2pt]
{\small
\begin{tabularx}{\textwidth}{>{\raggedright}p{2.1cm}>{\raggedright}p{3.6cm}>{\raggedright}X>{\raggedright\arraybackslash}X}
\toprule
Concern & Evidence (measured) & Limitation found & Improvement proposed\\
\midrule
Data and storage & 2,772 of 3,000 replies kept; 162 rejected for invented slots; labelling took 76\,min with 3 workers &
Single-replica SeaweedFS on a local volume; wrong-but-valid teacher labels pass the cleaning & Replicated
object store (3 volume servers); a second teacher pass or self-consistency vote to flag suspicious labels\\
Model quality & full frame 45\% (teacher 54\%), JSON validity 100\% on validation & Weak on places (82\%) and
intent (86\%); copies teacher errors & Train on all 2,772 records; constrained decoding of the intent enum;
better teacher prompt\\
Latency and throughput & p50 1,300\,ms, p95 1,700\,ms, 3.9 req/s at 4 users & One request at a time per pod, CPU-bound
greedy decoding & Batch concurrent requests, int8 quantisation or ONNX Runtime, KV-cache reuse of the
fixed prompt prefix\\
Resource utilisation & pods at 980--998m of 1 CPU, 583--620\,Mi of 1\,Gi & CPU saturates at 4 users; float16$\to$float32
conversion raises load-time memory to $\approx$790\,MiB & Store float32 (or load with \code{low\_cpu\_mem\_usage}) to cut
the peak; raise HPA max and lower the target so it scales before
saturation\\
Reliability and recovery & 0 failures in rolling update (1,314 req) and pod deletion (774 req); cold start 4\,s
& One laptop, one storage pod: no machine-level fault tolerance & Multi-node cluster across hosts, PodDisruptionBudget,
anti-affinity, replicated storage\\
Security and cost & gitleaks: 0 findings; 0.561 INR per 1,000 requests & Plain-HTTP ingress, one S3 admin key for
all workloads, billing on limits while idle & TLS on the Ingress, per-workload read-only S3 keys, NetworkPolicies,
scale to 1 replica when idle (already via HPA)\\
\bottomrule
\end{tabularx}}

\section*{References}
{\small Kim \& Rush, \emph{Sequence-Level Knowledge Distillation}, EMNLP 2016.
Hinton, Vinyals \& Dean, \emph{Distilling the Knowledge in a Neural Network}, 2015.
Bastianelli et al., \emph{SLURP}, EMNLP 2020 (CC BY 4.0).
Biderman et al., \emph{Pythia}, 2023 / EleutherAI Pythia-70M model card.
Documentation: Hugging Face Transformers \& TRL, k3d, K3s, SeaweedFS, Kubernetes (probes, HPA, Deployments,
Indexed Jobs), AWS VPC Network ACLs, Google SRE book (SLOs).}
\end{document}
